% Phụ lục A: Mã nguồn hệ thống BrainSeg-XAI

\chapter{Mã nguồn hệ thống BrainSeg-XAI}

\label{AppendixA}

%----------------------------------------------------------------------------------------

Phụ lục này trình bày toàn bộ mã nguồn của các module cốt lõi trong hệ thống BrainSeg-XAI.
Mã nguồn được tổ chức theo cấu trúc thư mục thực tế của dự án.

%----------------------------------------------------------------------------------------
\section{Backend -- Ứng dụng FastAPI chính (\texttt{app/main.py})}
%----------------------------------------------------------------------------------------

\begin{lstlisting}[style=codePython, caption={FastAPI -- Khoi tao app va CORS middleware (app/main.py)}]
from fastapi import FastAPI
from fastapi.middleware.cors import CORSMiddleware
from app.api.predict import router as predict_router

app = FastAPI(
    title="Brain Tumor AI API",
    description="He thong AI phan tich va danh gia khoi u nao tu anh MRI",
    version="1.0.0",
)

app.add_middleware(
    CORSMiddleware,
    allow_origins=["*"],
    allow_credentials=True,
    allow_methods=["*"],
    allow_headers=["*"],
)

app.include_router(predict_router, prefix="/api/v1")

@app.get("/")
async def root():
    return {
        "name": "Brain Tumor AI",
        "version": "1.0.0",
        "docs": "/docs",
    }
\end{lstlisting}

%----------------------------------------------------------------------------------------
\section{Backend -- Endpoint dự đoán (\texttt{app/api/predict.py})}
%----------------------------------------------------------------------------------------

\begin{lstlisting}[style=codePython, caption={REST API endpoint nhan anh MRI va tra ket qua phan tich (app/api/predict.py)}]
from fastapi import APIRouter, UploadFile, File, HTTPException
from app.services.inference import predict
from app.utils.image_utils import load_image_from_bytes

router = APIRouter()

@router.post("/predict")
async def predict_endpoint(file: UploadFile = File(...)):
    if not file.content_type.startswith("image/"):
        raise HTTPException(
            status_code=400,
            detail="Chi chap nhan file anh.")

    raw = await file.read()
    if len(raw) > 20 * 1024 * 1024:
        raise HTTPException(
            status_code=413,
            detail="File anh qua lon (toi da 20MB).")

    try:
        img_rgb = load_image_from_bytes(raw)
    except ValueError as e:
        raise HTTPException(status_code=422, detail=str(e))

    try:
        result = predict(img_rgb)
    except Exception as e:
        raise HTTPException(
            status_code=500,
            detail=f"Loi xu ly mo hinh: {str(e)}")

    return result

@router.get("/health")
async def health():
    return {"status": "ok",
            "message": "Brain Tumor AI API dang hoat dong"}
\end{lstlisting}

%----------------------------------------------------------------------------------------
\section{Backend -- Tiện ích xử lý ảnh (\texttt{app/utils/image\_utils.py})}
%----------------------------------------------------------------------------------------

\begin{lstlisting}[style=codePython, caption={Tien xu ly anh MRI: doc, resize 256x256, chuan hoa ImageNet (app/utils/image\_utils.py)}]
import cv2
import numpy as np
from PIL import Image
import io
import base64

TARGET_SIZE = (256, 256)
MEAN = [0.485, 0.456, 0.406]   # ImageNet mean
STD  = [0.229, 0.224, 0.225]   # ImageNet std

def load_image_from_bytes(raw: bytes) -> np.ndarray:
    arr = np.frombuffer(raw, dtype=np.uint8)
    img = cv2.imdecode(arr, cv2.IMREAD_COLOR)
    if img is None:
        raise ValueError(
            "Khong the doc anh. Vui long thu lai voi file anh hop le.")
    return cv2.cvtColor(img, cv2.COLOR_BGR2RGB)

def preprocess(img_rgb: np.ndarray):
    """Tra ve (tensor NCHW float32, original_rgb resized 256x256)"""
    import torch
    resized = cv2.resize(img_rgb, TARGET_SIZE,
                         interpolation=cv2.INTER_LINEAR)
    norm = resized.astype(np.float32) / 255.0
    for c in range(3):
        norm[:, :, c] = (norm[:, :, c] - MEAN[c]) / STD[c]
    tensor = torch.from_numpy(
        norm.transpose(2, 0, 1)).unsqueeze(0)  # 1x3xHxW
    return tensor, resized

def mask_to_binary(logits_tensor) -> np.ndarray:
    import torch
    prob = torch.sigmoid(logits_tensor).squeeze().cpu().numpy()
    return (prob > 0.5).astype(np.uint8)

def ndarray_to_b64(arr: np.ndarray, ext: str = "png") -> str:
    pil = Image.fromarray(arr.astype(np.uint8))
    buf = io.BytesIO()
    pil.save(buf, format=ext.upper())
    return base64.b64encode(buf.getvalue()).decode()

def apply_colormap_heatmap(gray: np.ndarray) -> np.ndarray:
    """gray: HxW float [0,1] -> RGB uint8 (JET colormap)"""
    gray_u8 = (gray * 255).clip(0, 255).astype(np.uint8)
    colored = cv2.applyColorMap(gray_u8, cv2.COLORMAP_JET)
    return cv2.cvtColor(colored, cv2.COLOR_BGR2RGB)

def overlay_mask(img_rgb: np.ndarray, mask: np.ndarray,
                 alpha: float = 0.45) -> np.ndarray:
    overlay = img_rgb.copy()
    overlay[mask == 1] = [255, 80, 80]
    return cv2.addWeighted(img_rgb, 1 - alpha, overlay, alpha, 0)

def overlay_heatmap(img_rgb: np.ndarray, heatmap_rgb: np.ndarray,
                    alpha: float = 0.5) -> np.ndarray:
    return cv2.addWeighted(img_rgb, 1 - alpha, heatmap_rgb, alpha, 0)
\end{lstlisting}

%----------------------------------------------------------------------------------------
\section{Backend -- Định nghĩa mô hình (\texttt{app/models/unet.py})}
%----------------------------------------------------------------------------------------

\begin{lstlisting}[style=codePython, caption={Vanilla UNet va SMP wrapper -- ho tro U-Net va U-Net++ (app/models/unet.py)}]
import torch
import torch.nn as nn
import torch.nn.functional as F

# --- 1. Vanilla U-Net (backward-compat voi weight cu) ---------------
class DoubleConv(nn.Module):
    def __init__(self, in_ch, out_ch):
        super().__init__()
        self.conv = nn.Sequential(
            nn.Conv2d(in_ch, out_ch, 3, padding=1, bias=False),
            nn.BatchNorm2d(out_ch), nn.ReLU(inplace=True),
            nn.Conv2d(out_ch, out_ch, 3, padding=1, bias=False),
            nn.BatchNorm2d(out_ch), nn.ReLU(inplace=True),
        )
    def forward(self, x): return self.conv(x)

class Down(nn.Module):
    def __init__(self, ic, oc):
        super().__init__()
        self.pool_conv = nn.Sequential(
            nn.MaxPool2d(2), DoubleConv(ic, oc))
    def forward(self, x): return self.pool_conv(x)

class Up(nn.Module):
    def __init__(self, ic, oc):
        super().__init__()
        self.up   = nn.ConvTranspose2d(ic, ic // 2, 2, stride=2)
        self.conv = DoubleConv(ic, oc)
    def forward(self, x1, x2):
        x1 = self.up(x1)
        dy = x2.size(2) - x1.size(2)
        dx = x2.size(3) - x1.size(3)
        x1 = F.pad(x1, [dx//2, dx-dx//2, dy//2, dy-dy//2])
        return self.conv(torch.cat([x2, x1], 1))

class UNet(nn.Module):
    def __init__(self, n_channels=3, n_classes=1,
                 features=[64, 128, 256, 512]):
        super().__init__()
        f = features
        self.inc   = DoubleConv(n_channels, f[0])
        self.down1 = Down(f[0], f[1])
        self.down2 = Down(f[1], f[2])
        self.down3 = Down(f[2], f[3])
        self.down4 = Down(f[3], f[3]*2)
        self.up1   = Up(f[3]*2, f[3])
        self.up2   = Up(f[3],   f[2])
        self.up3   = Up(f[2],   f[1])
        self.up4   = Up(f[1],   f[0])
        self.outc  = nn.Conv2d(f[0], n_classes, 1)

    def forward(self, x):
        x1 = self.inc(x)
        x2 = self.down1(x1)
        x3 = self.down2(x2)
        x4 = self.down3(x3)
        x5 = self.down4(x4)
        x  = self.up1(x5, x4)
        x  = self.up2(x,  x3)
        x  = self.up3(x,  x2)
        x  = self.up4(x,  x1)
        return self.outc(x)

# --- 2. SMP wrapper --------------------------------------------------
def build_smp_model(arch="unetplusplus",
                    encoder="efficientnet-b4",
                    encoder_weights="imagenet",
                    n_classes=1):
    """
    arch    : 'unetplusplus' | 'unet' | 'deeplabv3plus' | 'fpn' | 'manet'
    encoder : 'efficientnet-b4' | 'resnet50' | 'resnet34' | ...
    """
    try:
        import segmentation_models_pytorch as smp
    except ImportError:
        raise ImportError(
            "segmentation-models-pytorch chua duoc cai.\n"
            "Chay: pip install segmentation-models-pytorch timm")
    return smp.create_model(
        arch,
        encoder_name=encoder,
        encoder_weights=encoder_weights,
        in_channels=3,
        classes=n_classes,
        activation=None,
    )

# --- 3. Auto-loader --------------------------------------------------
def get_model(arch="unetplusplus", encoder="efficientnet-b4",
              n_classes=1, pretrained=True):
    try:
        model = build_smp_model(
            arch=arch, encoder=encoder,
            encoder_weights="imagenet" if pretrained else None,
            n_classes=n_classes)
        print(f"[Model] Dung SMP -- arch={arch}, encoder={encoder}")
        return model, "smp"
    except Exception as e:
        print(f"[Model] SMP khong kha dung ({e}), fallback UNet.")
        return UNet(n_channels=3, n_classes=n_classes), "vanilla"
\end{lstlisting}

%----------------------------------------------------------------------------------------
\section{Backend -- Dịch vụ suy luận (\texttt{app/services/inference.py})}
%----------------------------------------------------------------------------------------

\begin{lstlisting}[style=codePython, caption={Lazy model loading, auto-detect checkpoint va dieu phoi prediction pipeline (app/services/inference.py)}]
import torch
import numpy as np
from pathlib import Path

from app.models.unet import get_model, UNet
from app.utils.image_utils import (
    preprocess, mask_to_binary,
    apply_colormap_heatmap, overlay_mask, overlay_heatmap,
    ndarray_to_b64,
)
from app.xai.gradcam import GradCAM, get_target_layer
from app.services.feature_extraction import extract_features
from app.rules.medical_rules import evaluate_risk

WEIGHTS_PATH = (Path(__file__).parent.parent.parent
                / "weights" / "unet_best.pth")
DEVICE = torch.device("cuda" if torch.cuda.is_available() else "cpu")

_model      = None
_model_type = "vanilla"
_gradcam    = None

def convert_numpy(obj):
    if isinstance(obj, (np.bool_, bool)):    return bool(obj)
    if isinstance(obj, np.integer):          return int(obj)
    if isinstance(obj, np.floating):         return float(obj)
    if isinstance(obj, np.ndarray):          return obj.tolist()
    if isinstance(obj, dict):
        return {k: convert_numpy(v) for k, v in obj.items()}
    if isinstance(obj, (list, tuple)):
        return [convert_numpy(v) for v in obj]
    return obj

def _load_model():
    global _model, _model_type, _gradcam
    if _model is not None:
        return

    if not WEIGHTS_PATH.exists():
        print("[WARN] Weights not found -- demo mode (random weights).")
        model, mtype = get_model(pretrained=False)
        model.to(DEVICE).eval()
        _model      = model
        _model_type = mtype
        _gradcam    = GradCAM(model, get_target_layer(model, mtype))
        return

    ckpt = torch.load(WEIGHTS_PATH, map_location=DEVICE)

    arch    = (ckpt.get("arch",    "unetplusplus")
               if isinstance(ckpt, dict) else "unetplusplus")
    encoder = (ckpt.get("encoder", "efficientnet-b4")
               if isinstance(ckpt, dict) else "efficientnet-b4")

    if isinstance(ckpt, dict) and "model_state_dict" in ckpt:
        sd = ckpt["model_state_dict"]
    elif isinstance(ckpt, dict) and "state_dict" in ckpt:
        sd = ckpt["state_dict"]
    elif isinstance(ckpt, dict) and any(
        k.startswith("encoder") or k.startswith("decoder")
        for k in ckpt.keys()
    ):
        sd = ckpt
        arch, encoder = "unetplusplus", "efficientnet-b4"
    else:
        sd = ckpt

    keys = list(sd.keys()) if isinstance(sd, dict) else []
    is_vanilla = any(
        k.startswith("inc.") or k.startswith("down")
        or k.startswith("up")
        for k in keys)

    if is_vanilla:
        model, mtype = UNet(n_channels=3, n_classes=1), "vanilla"
        print(f"[Model] Load vanilla UNet tu {WEIGHTS_PATH}")
    else:
        try:
            model, mtype = get_model(
                arch=arch, encoder=encoder, pretrained=False)
            print(f"[Model] Load SMP {arch}/{encoder}")
        except Exception as e:
            print(f"[Model] SMP fail ({e}), thu vanilla UNet.")
            model, mtype = UNet(n_channels=3, n_classes=1), "vanilla"

    model.load_state_dict(sd, strict=False)
    model.to(DEVICE).eval()
    _model      = model
    _model_type = mtype
    _gradcam    = GradCAM(model, get_target_layer(model, mtype))
    print(f"[Model] Ready -- type={mtype}, device={DEVICE}")

def predict(img_rgb: np.ndarray) -> dict:
    _load_model()
    tensor, resized = preprocess(img_rgb)
    tensor = tensor.to(DEVICE)

    with torch.no_grad():
        logits = _model(tensor)

    mask = mask_to_binary(logits)

    try:
        cam = _gradcam.generate(tensor.clone())
    except Exception as e:
        print(f"[GradCAM] Error: {e}")
        cam = np.zeros(resized.shape[:2], dtype=np.float32)

    heatmap_rgb  = apply_colormap_heatmap(cam)
    mask_overlay = overlay_mask(resized, mask)
    cam_overlay  = overlay_heatmap(resized, heatmap_rgb)

    mask_vis = np.zeros((*mask.shape, 3), dtype=np.uint8)
    mask_vis[mask == 1] = [0, 255, 136]

    features = extract_features(mask, mask.shape)
    risk     = evaluate_risk(features)

    return convert_numpy({
        "original_b64":    ndarray_to_b64(resized),
        "mask_b64":        ndarray_to_b64(mask_vis),
        "overlay_b64":     ndarray_to_b64(mask_overlay),
        "heatmap_b64":     ndarray_to_b64(heatmap_rgb),
        "cam_overlay_b64": ndarray_to_b64(cam_overlay),
        "features": features,
        "risk": {
            "risk_level":      risk.risk_level,
            "risk_score":      risk.risk_score,
            "severity":        risk.severity,
            "risk_color":      risk.risk_color,
            "fired_rules":     risk.fired_rules,
            "recommendations": risk.recommendations,
            "explanation":     risk.explanation,
        },
    })
\end{lstlisting}

%----------------------------------------------------------------------------------------
\section{Backend -- Module Grad-CAM (\texttt{app/xai/gradcam.py})}
%----------------------------------------------------------------------------------------

\begin{lstlisting}[style=codePython, caption={Trien khai Grad-CAM voi PyTorch forward/backward hooks (app/xai/gradcam.py)}]
import numpy as np
import torch
import torch.nn.functional as F

class GradCAM:
    def __init__(self, model: torch.nn.Module,
                 target_layer: torch.nn.Module):
        self.model        = model
        self.target_layer = target_layer
        self._features    = None
        self._grads       = None
        self._hooks       = []
        self._register()

    def _register(self):
        self._hooks.append(
            self.target_layer.register_forward_hook(
                self._save_features))
        self._hooks.append(
            self.target_layer.register_full_backward_hook(
                self._save_grads))

    def _save_features(self, m, inp, out):
        self._features = out.detach()

    def _save_grads(self, m, gi, go):
        self._grads = go[0].detach()

    def remove_hooks(self):
        for h in self._hooks:
            h.remove()

    def generate(self, input_tensor: torch.Tensor) -> np.ndarray:
        self.model.eval()
        input_tensor = input_tensor.detach().requires_grad_(True)

        output = self.model(input_tensor)
        score  = torch.sigmoid(output).sum()   # tong sigmoid cho phan doan nhi phan
        self.model.zero_grad()
        score.backward()

        if self._grads is None or self._features is None:
            return np.zeros(input_tensor.shape[-2:], dtype=np.float32)

        # Global Average Pooling tren gradient
        weights = self._grads.mean(dim=(2, 3), keepdim=True)
        cam = F.relu(
            (weights * self._features).sum(dim=1, keepdim=True))

        cam = F.interpolate(
            cam,
            size=input_tensor.shape[-2:],
            mode="bilinear",
            align_corners=False)
        cam = cam.squeeze().cpu().numpy()

        lo, hi = cam.min(), cam.max()
        if hi - lo > 1e-8:
            cam = (cam - lo) / (hi - lo)  # chuan hoa min-max
        else:
            cam = np.zeros_like(cam)
        return cam.astype(np.float32)

def get_target_layer(model, model_type: str = "vanilla"):
    """Chon layer de dang ky hook Grad-CAM theo loai model."""
    if model_type == "smp":
        try:
            return model.segmentation_head[0]
        except Exception:
            pass
        try:
            blocks = list(model.decoder.blocks)
            return blocks[-1].conv2[0]
        except Exception:
            pass
        convs = [m for m in model.modules()
                 if isinstance(m, torch.nn.Conv2d)]
        return convs[-2] if len(convs) >= 2 else convs[-1]
    else:
        # Vanilla UNet: layer truoc lop dau ra cuoi cung
        return model.up4.conv.conv[-3]
\end{lstlisting}

%----------------------------------------------------------------------------------------
\section[Backend -- Trích xuất đặc trưng]{Backend -- Trích xuất đặc trưng (\texttt{feature\_extraction.py})}
%----------------------------------------------------------------------------------------

\begin{lstlisting}[style=codePython, caption={Tinh toan dac trung hinh thai hoc tu mat na nhi phan (app/services/feature\_extraction.py)}]
import numpy as np
from scipy import ndimage
from skimage import measure

def extract_features(mask: np.ndarray, img_shape: tuple) -> dict:
    """
    mask      : binary HxW uint8
    img_shape : (H, W)
    Returns   : dict cac dac trung y khoa
    """
    H, W = img_shape
    total_pixels = H * W

    labeled, num_tumors = ndimage.label(mask)
    tumor_pixels    = int(mask.sum())
    occupancy_ratio = round(tumor_pixels / total_pixels * 100, 2)

    if tumor_pixels == 0:
        return {
            "tumor_detected":    False,
            "tumor_area_px":     0,
            "tumor_area_cm2":    0.0,
            "occupancy_ratio":   0.0,
            "num_regions":       0,
            "shape_irregularity": 0.0,
            "compactness":       0.0,
            "boundary_complexity": 0.0,
            "midline_shift":     False,
            "location":          "Khong phat hien khoi u",
            "centroid_x":        None,
            "centroid_y":        None,
        }

    # Chuyen doi px -> cm2 (MRI slice ~ 24cm x 24cm tai 256x256)
    px_to_cm = (24.0 / H) ** 2
    tumor_area_cm2 = round(tumor_pixels * px_to_cm, 2)

    # Dac trung hinh dang qua regionprops cua vung lon nhat
    props_list = measure.regionprops(labeled)
    props = max(props_list, key=lambda p: p.area)

    perimeter = props.perimeter if props.perimeter > 0 else 1
    area_p    = props.area
    compactness = round((4 * np.pi * area_p) / (perimeter ** 2), 4)
    shape_irregularity  = round(1.0 - compactness, 4)
    boundary_complexity = round(
        perimeter / (np.sqrt(area_p) + 1e-6), 4)

    cy, cx    = props.centroid
    cx_norm   = cx / W
    cy_norm   = cy / H

    h_loc = ("ben trai" if cx_norm < 0.45
             else "ben phai" if cx_norm > 0.55
             else "trung tam")
    v_loc = ("thuy tran" if cy_norm < 0.33
             else "thuy dinh/thai duong" if cy_norm < 0.66
             else "thuy cham/tieu nao")
    location = f"{v_loc}, {h_loc}"

    midline_shift = abs(cx_norm - 0.5) > 0.10

    return {
        "tumor_detected":    True,
        "tumor_area_px":     tumor_pixels,
        "tumor_area_cm2":    tumor_area_cm2,
        "occupancy_ratio":   occupancy_ratio,
        "num_regions":       num_tumors,
        "shape_irregularity": shape_irregularity,
        "compactness":       compactness,
        "boundary_complexity": boundary_complexity,
        "midline_shift":     midline_shift,
        "location":          location,
        "centroid_x":        round(float(cx), 1),
        "centroid_y":        round(float(cy), 1),
    }
\end{lstlisting}

%----------------------------------------------------------------------------------------
\section{Backend -- Bộ máy luật y khoa (\texttt{app/rules/medical\_rules.py})}
%----------------------------------------------------------------------------------------

\begin{lstlisting}[style=codePython, caption={Rule-based risk engine -- 6 quy tac, additive scoring, 4 muc rui ro (app/rules/medical\_rules.py)}]
from dataclasses import dataclass, field
from typing import List

@dataclass
class RuleResult:
    risk_level:      str   # "Thap" | "Trung binh" | "Cao" | "Rat cao"
    risk_score:      int   # 0-100
    severity:        str   # "Binh thuong" | "Nhe" | ... | "Nguy kich"
    risk_color:      str   # hex color cho UI
    fired_rules:     List[str] = field(default_factory=list)
    recommendations: List[str] = field(default_factory=list)
    explanation:     str = ""

def evaluate_risk(features: dict) -> RuleResult:
    if not features.get("tumor_detected", False):
        return RuleResult(
            risk_level="Thap", risk_score=0,
            severity="Binh thuong", risk_color="#00ff88",
            fired_rules=["Khong phat hien khoi u trong anh MRI"],
            recommendations=["Tiep tuc theo doi dinh ky",
                             "Chup MRI kiem tra sau 12 thang"],
            explanation="Khong phat hien bat ky vung bat thuong nao.")

    score = 0
    fired, recs = [], []

    occ   = features["occupancy_ratio"]
    irr   = features["shape_irregularity"]
    bnd   = features["boundary_complexity"]
    mid   = features["midline_shift"]
    n_reg = features["num_regions"]
    area  = features["tumor_area_cm2"]

    # Rule 1: Occupancy ratio
    if occ > 20:
        score += 30
        fired.append(f"Ty le chiem vung nao rat cao ({occ}% > 20%)")
        recs.append("Can can thiep phau thuat khan cap")
    elif occ > 10:
        score += 20
        fired.append(f"Ty le chiem vung nao cao ({occ}% > 10%)")
        recs.append("Hoi chan chuyen khoa than kinh ngay")
    elif occ > 5:
        score += 10
        fired.append(f"Ty le chiem vung nao o muc trung binh ({occ}%)")
        recs.append("Theo doi sat va chup MRI dinh ky 3 thang/lan")
    else:
        fired.append(f"Ty le chiem vung nao thap ({occ}%)")

    # Rule 2: Shape irregularity
    if irr > 0.7:
        score += 25
        fired.append(
            f"Hinh dang khoi u rat bat doi xung (irr={irr:.2f})")
        recs.append("Sinh thiet de xac dinh tinh ac tinh")
    elif irr > 0.5:
        score += 15
        fired.append(
            f"Hinh dang khoi u bat doi xung vua (irr={irr:.2f})")
    else:
        fired.append(
            f"Hinh dang khoi u tuong doi deu dan (irr={irr:.2f})")

    # Rule 3: Boundary complexity
    if bnd > 15:
        score += 20
        fired.append(
            f"Duong bien khoi u rat phuc tap (complexity={bnd:.1f})")
        recs.append("Chi dinh PET-CT de danh gia do xam lan")
    elif bnd > 10:
        score += 10
        fired.append(
            f"Duong bien khoi u phuc tap vua (complexity={bnd:.1f})")
    else:
        fired.append(
            f"Duong bien khoi u tuong doi don gian (complexity={bnd:.1f})")

    # Rule 4: Midline shift
    if mid:
        score += 20
        fired.append(
            "Co dau hieu dich chuyen duong giua nao (Midline Shift)")
        recs.append(
            "Theo doi ap luc noi so, xem xet can thiep giai ap")

    # Rule 5: Multiple regions
    if n_reg > 3:
        score += 15
        fired.append(
            f"Phat hien nhieu vung khoi u ({n_reg} vung) -- nghi di can")
        recs.append("Chup MRI toan than, tim kiem o di can nguyen phat")
    elif n_reg > 1:
        score += 8
        fired.append(f"Phat hien {n_reg} vung khoi u rieng biet")

    # Rule 6: Absolute size
    if area > 10:
        score += 15
        fired.append(f"Dien tich khoi u rat lon ({area} cm2)")
    elif area > 4:
        score += 8
        fired.append(f"Dien tich khoi u trung binh ({area} cm2)")
    else:
        fired.append(f"Dien tich khoi u nho ({area} cm2)")

    score = min(score, 100)

    if score >= 70:
        risk_level, severity, risk_color = (
            "Rat cao", "Nguy kich", "#ff2d55")
        recs.insert(0,
            "KHAN CAP: Can nhap vien va can thiep ngay lap tuc")
    elif score >= 50:
        risk_level, severity, risk_color = (
            "Cao", "Nghiem trong", "#ff6b35")
        recs.insert(0,
            "Can dieu tri chuyen khoa trong vong 48 gio")
    elif score >= 30:
        risk_level, severity, risk_color = (
            "Trung binh", "Trung binh", "#ffd60a")
        if not recs:
            recs.append("Theo doi va dieu tri theo phac do cua bac si")
    else:
        risk_level, severity, risk_color = (
            "Thap", "Nhe", "#00c896")
        if not recs:
            recs.append("Theo doi dinh ky 6 thang/lan")

    explanation = (
        f"He thong phan tich {len(fired)} dac trung y khoa va kich hoat "
        f"{sum(1 for r in fired if any(k in r for k in ['cao', 'phuc tap', 'bat doi', 'dich chuyen', 'di can', 'lon']))} "
        f"co nguy co. Diem nguy co tong hop: {score}/100. "
        f"Phan loai: {severity}."
    )

    return RuleResult(
        risk_level=risk_level, risk_score=score,
        severity=severity, risk_color=risk_color,
        fired_rules=fired,
        recommendations=list(dict.fromkeys(recs)),
        explanation=explanation,
    )
\end{lstlisting}

%----------------------------------------------------------------------------------------
\section[Frontend -- Thành phần xem ảnh]{Frontend -- Thành phần xem ảnh (\texttt{ImageViewer.jsx})}
%----------------------------------------------------------------------------------------

\begin{lstlisting}[style=codePython, caption={React component ImageViewer -- hien thi 5 che do xem anh (ImageViewer.jsx)}]
import { useState } from 'react';

const VIEWS = [
  { key: 'original_b64', label: 'Anh Goc',      tag: 'MRI RAW'       },
  { key: 'overlay_b64',  label: 'Mask Overlay',  tag: 'SEGMENTATION'  },
  { key: 'mask_b64',     label: 'Segmentation',  tag: 'BINARY MASK'   },
  { key: 'cam_overlay_b64', label: 'XAI Overlay', tag: 'GRAD-CAM'    },
  { key: 'heatmap_b64',  label: 'Heatmap',       tag: 'ATTENTION MAP' },
];

export default function ImageViewer({ result }) {
  const [active, setActive] = useState(0);
  const current = VIEWS[active];
  const src = `data:image/png;base64,${result[current.key]}`;

  return (
    <div className="card animate-in">
      <div className="card-header">
        <span className="card-title">Truc quan hoa ket qua</span>
        <span className="card-tag">{current.tag}</span>
      </div>

      {/* Tab selector */}
      <div className="image-tabs">
        {VIEWS.map((v, i) => (
          <button
            key={v.key}
            className={`img-tab${active === i ? ' active' : ''}`}
            onClick={() => setActive(i)}
          >
            {v.label}
          </button>
        ))}
      </div>

      {/* Anh chinh + thumbnails */}
      <div className="image-display">
        <div className="img-main-wrapper">
          <img className="img-main" src={src} alt={current.label} />
          <div className="img-label">{current.label}</div>
        </div>
        <div className="img-thumbs">
          {VIEWS.map((v, i) => (
            <div key={v.key}>
              <div
                className={`img-thumb${active === i ? ' active' : ''}`}
                onClick={() => setActive(i)}
              >
                <img
                  src={`data:image/png;base64,${result[v.key]}`}
                  alt={v.label}
                />
              </div>
              <div className="thumb-label">
                {v.label.split(' ')[0]}
              </div>
            </div>
          ))}
        </div>
      </div>
    </div>
  );
}
\end{lstlisting}

%----------------------------------------------------------------------------------------
\section[Frontend -- Bảng đánh giá rủi ro]{Frontend -- Bảng đánh giá rủi ro (\texttt{RiskPanel.jsx})}
%----------------------------------------------------------------------------------------

\begin{lstlisting}[style=codePython, caption={React component RiskPanel -- dong ho rui ro SVG va ket qua phan tich (RiskPanel.jsx)}]
import { useEffect, useState } from 'react';

function RiskGauge({ score, color }) {
  const [displayed, setDisplayed] = useState(0);
  const R = 54;
  const CIRC = 2 * Math.PI * R;

  useEffect(() => {
    let start = null;
    const duration = 1200;
    const animate = (ts) => {
      if (!start) start = ts;
      const prog = Math.min((ts - start) / duration, 1);
      const ease = 1 - Math.pow(1 - prog, 3); // cubic ease-out
      setDisplayed(Math.round(score * ease));
      if (prog < 1) requestAnimationFrame(animate);
    };
    requestAnimationFrame(animate);
  }, [score]);

  const offset = CIRC * (1 - displayed / 100);
  return (
    <div className="risk-gauge">
      <svg width="140" height="140" viewBox="0 0 140 140">
        <circle cx="70" cy="70" r={R} fill="none"
          stroke="rgba(0,200,255,0.07)" strokeWidth="10" />
        <circle cx="70" cy="70" r={R} fill="none"
          stroke={color} strokeWidth="10"
          strokeLinecap="round"
          strokeDasharray={CIRC}
          strokeDashoffset={offset}
          style={{
            transform: 'rotate(-90deg)',
            transformOrigin: '70px 70px',
            filter: `drop-shadow(0 0 8px ${color})`,
          }}
        />
      </svg>
      <div className="risk-number">
        <span style={{ color }}>{displayed}</span>
        <small>/100</small>
      </div>
    </div>
  );
}

export default function RiskPanel({ risk, features }) {
  const color = risk.risk_color;
  const levelEmoji = {
    'Rat cao': 'Nguy kich', 'Cao': 'Nghiem trong',
    'Trung binh': 'Trung binh', 'Thap': 'Nhe',
  }[risk.risk_level] || '';

  return (
    <div style={{ display: 'flex', flexDirection: 'column', gap: 20 }}>

      {/* Diem nguy co */}
      <div className="card risk-card animate-in">
        <div className="card-header">
          <span className="card-title">Danh gia nguy co lam sang</span>
        </div>
        <div className="risk-score-display">
          <RiskGauge score={risk.risk_score} color={color} />
          <div className="risk-level-badge"
            style={{ color, border: `1px solid ${color}40` }}>
            Nguy co {risk.risk_level}
          </div>
          <div className="risk-severity">
            Muc do: <span style={{ color }}>{risk.severity}</span>
          </div>
        </div>
      </div>

      {/* Luat da kich hoat */}
      <div className="card animate-in">
        <div className="card-header">
          <span className="card-title">Luat y khoa kich hoat</span>
          <span className="card-tag">{risk.fired_rules.length} luat</span>
        </div>
        <div className="rules-list">
          {risk.fired_rules.map((rule, i) => (
            <div key={i} className="rule-item">
              <span className="rule-icon">
                {rule.includes('cao') ? 'WARN' : 'OK'}
              </span>
              <span>{rule}</span>
            </div>
          ))}
        </div>
      </div>

      {/* Khuyen nghi lam sang */}
      <div className="card animate-in">
        <div className="card-header">
          <span className="card-title">Khuyen nghi lam sang</span>
        </div>
        <div className="rec-list">
          {risk.recommendations.map((rec, i) => (
            <div key={i} className="rec-item">
              <span>-- {rec}</span>
            </div>
          ))}
        </div>
      </div>

    </div>
  );
}
\end{lstlisting}

%----------------------------------------------------------------------------------------
\section{Frontend -- Dịch vụ API client (\texttt{src/services/api.js})}
%----------------------------------------------------------------------------------------

\begin{lstlisting}[style=codePython, caption={Axios client goi REST API backend (src/services/api.js)}]
import axios from 'axios';

const API_BASE = process.env.REACT_APP_API_URL
                 || 'http://localhost:8000';

const api = axios.create({
  baseURL: `${API_BASE}/api/v1`,
  timeout: 120000,   // 2 phut timeout cho inference tren CPU
});

export const predictTumor = async (file, onUploadProgress) => {
  const formData = new FormData();
  formData.append('file', file);

  const response = await api.post('/predict', formData, {
    headers: { 'Content-Type': 'multipart/form-data' },
    onUploadProgress,
  });

  return response.data;
};

export const checkHealth = async () => {
  const response = await api.get('/health');
  return response.data;
};
\end{lstlisting}

%----------------------------------------------------------------------------------------
\section{Docker Compose -- Cấu hình triển khai (\texttt{docker-compose.yml})}
%----------------------------------------------------------------------------------------

\begin{lstlisting}[style=plaintext, caption={Cau hinh Docker Compose cho trien khai production (docker-compose.yml)}]
version: '3.8'

services:
  backend:
    build: ./backend
    ports:
      - "8000:8000"
    volumes:
      - ./backend/weights:/app/weights
    environment:
      - PYTHONUNBUFFERED=1
    restart: unless-stopped

  frontend:
    build: ./frontend
    ports:
      - "3000:80"
    depends_on:
      - backend
    environment:
      - REACT_APP_API_URL=http://backend:8000
    restart: unless-stopped
\end{lstlisting}

%----------------------------------------------------------------------------------------
\section{Đoạn trích minh họa -- Nhận dạng và lưu checkpoint}
%----------------------------------------------------------------------------------------

Mục này trích xuất hai đoạn mã ngắn minh họa cơ chế nhận dạng loại checkpoint và định dạng lưu checkpoint chuẩn được đề cập trong Chương~\ref{Chapter3}.

\subsection*{A.13.1 Nhận dạng loại checkpoint theo tên khóa state\_dict}

\begin{lstlisting}[style=codePython, caption={Nhan dang loai checkpoint theo ten khoa state\_dict (app/services/inference.py)}]
# doc state_dict tu file .pth
sd = checkpoint if isinstance(checkpoint, dict) \
     and "model_state_dict" not in checkpoint \
     else checkpoint.get("model_state_dict", checkpoint)

keys = list(sd.keys())
is_vanilla = any(
    k.startswith("inc.") or k.startswith("down") or k.startswith("up")
    for k in keys
)

if is_vanilla:
    model = UNet(n_channels=3, n_classes=1)
else:
    arch    = checkpoint.get("arch",    "unetplusplus")
    encoder = checkpoint.get("encoder", "efficientnet-b4")
    model   = build_smp_model(arch, encoder)

model.load_state_dict(sd)
model.eval()
\end{lstlisting}

\subsection*{A.13.2 Lưu checkpoint kèm metadata kiến trúc}

\begin{lstlisting}[style=codePython, caption={Luu checkpoint kem metadata kien truc (dung khi ket thuc huan luyen)}]
# Luu sau khi huan luyen de he thong tu dong nhan dang khi tai lai
torch.save(
    {
        "arch":             "unetplusplus",
        "encoder":          "efficientnet-b4",
        "model_state_dict": model.state_dict(),
        "epoch":            best_epoch,
        "val_dice":         best_dice,
    },
    "backend/weights/unet_best.pth"
)
\end{lstlisting}
